% (C) 2016 Jean Nassar. Some Rights Reserved
% Except where otherwise noted, this work is licensed under the Creative Commons Attribution-ShareAlike License, version 4
% Title page
\renewcommand{\maketitlehooka}{
  \begin{center}
    \noindent

    \begin{figure}[h]
      \centering
      \includegraphics[width=0.3\textwidth]{kyoto_u_logo}
    \end{figure}
    \textsc{\LARGE Kyoto University}
    
    \textsc{\Large Graduate School of Engineering}

    \textsc{\large Mechanical Engineering Science}
  \end{center}
  \vfill\vfill
}
\renewcommand{\maketitlehookb}{
  \vfill\vfill
  \begin{center}
  {\large A thesis presented for the degree of}

  \textsc{\large Master of Science}
  \end{center}
  \vfill\vfill
}
\renewcommand{\maketitlehookc}{
  %\vfill
}
%\renewcommand{\maketitlehookd}{
  %\begin{center}
    %\begin{figure}[h]
      %\centering
      %\includegraphics{by-sa-big}
    %\end{figure}
  %\end{center}
%}

\title{\mytitle}
% \author{Jean \textsc{Nassar}\thanks{Supervisor: Dr.~Fumitoshi Matsuno}}
\author{Jean \textsc{Nassar} \\\emph{Author} \and Dr.~Fumitoshi \textsc{Matsuno}\\\emph{Supervisor}}
\date{\today}

\begin{titlingpage}
\maketitle
\thispagestyle{empty}

% Copyright page
\clearpage
\null\vfill
\noindent \emph{\mytitle}

\noindent \emph{\mytitlejp}
\vspace{1em}

\noindent All the source code for this project is accessible at the author's public Github profile at \url{https://github.com/masasin/spirit}.
\vspace{1em}

\noindent This document was typeset using the \XeTeX\ typesetting system created by the Non-Roman Script Initiative and the \textsf{memoir} class created by Peter Wilson.
The text is set in \mytextsize\ in Linux~Libertine~O.
The other fonts used include \textsf{Linux~Biolinum~O}, \texttt{DejaVu~Sans~Code}, and IPA~Mincho~(明朝).
\vspace{1em}

\noindent The PDF version of this thesis is fully hyperlinked and navigable.
\vspace{1em}

\noindent \includegraphics{by-sa-big}

\noindent
% \includegraphics{by-sa}
Copyright \copyright 2017 by Jean Nassar.

\noindent This thesis is licensed under the Creative Commons Attribution-ShareAlike 4.0 International License.
To view a copy of this license, visit:\\
\url{http://creativecommons.org/licenses/by-sa/4.0/}

\noindent The \acrshort{spirit} software, as well as the code used for data collection and analysis, are licensed under the MIT License.

\end{titlingpage}


% Dedication
%\clearpage
\thispagestyle{empty}
\setlength{\epigraphwidth}{0.7\textwidth}
\null\vfill
\epigraph{
  because [\ldots] if you don't know where you are, then you don't know where you're going.
  And if you don't know where you're going, you're probably going wrong.}
  {\emph{I Shall Wear Midnight}\\
   Terry \textsc{Pratchett}}
\vfill\vfill\null
  %\null\vspace{\stretch{1}}
    %\noindent
    %because [\ldots] if you don't know where you are, then you don't know where you're going.
    %And if you don't know where you're going, you're probably going wrong.

    %\begin{flushright}
      %\emph{I Shall Wear Midnight}\\
      %\textsc{Terry Pratchett}
    %\end{flushright}
  %\vspace{\stretch{2}}\null

% Abstract
\clearpage
\thispagestyle{empty}
\null\vfill
\renewcommand{\abstractname}{Summary}
\begin{abstract}
\gls{spirit} is a novel robotics teleoperation system which overlays the current position and orientation of a vehicle onto previously acquired images.
This research focuses on developing a \gls{spirit}-based user interface for aerial robots.
The proposed method combines \gls{fov} information with state estimation, and selects a suitable image to use as a background.
It works even in a low-throughput environment where video data might suffer, and only requires one camera.

Experiments were conducted to study the efficacy at increasing accuracy while decreasing time and wasted movements.
Participants were tasked with flying a drone to a location above a known target while using the proposed system, and comparing the results with those obtained by flying with a slowed-down onboard camera.

\gls{spirit} significantly improved accuracy from 0.666\,m to 0.401\,m, an improvement of almost 40\% ($\Delta$\sym{mean}=0.266\,m, \sym{effect}=1.053, \acrshort{ci}=98.1\%), and reduced the overall workload by 37.5\% ($p$=0.02408, \sym{effect}=$-0.978$), especially in the physical, temporal, and performance metrics.
It also significantly improved the users' control and awareness of absolute and relative positioning.
Their self-assessed ability to stay above the target increased by 86.9\% from an average rating of $2.556/7$ to $4.778/7$ ($p$=0.00126, \sym{effect}=2.511).

There were non-significant increases in time and path lengths, but improvements over subsequent runs indicate that these are simply an issue of familiarity with the system.
Further testing is needed to verify the effect of \gls{spirit} on completion time and path length, and its efficacy in various environments.

Overall, the tests were successful, and all participants gave praise for the system, as well as valuable feedback.

\end{abstract}
\vfill\vfill\null


% Acknowledgements
\clearpage
\thispagestyle{empty}
\null\vfill
\renewcommand{\abstractname}{Acknowledgements}
\begin{abstract}
  I would first like to thank my thesis advisor, Professor Fumitoshi Matsuno for his guidance during my past several years here.
  He never gave up on me even when I was discouraged.
  I would also like to thank the staff and students at Kyoto University for their valuable advice and companionship.
  From building robots to intramurals, from dinner invitations to hugs, hiking and flying, I had an unforgettable experience.
  Thank you especially to the participants in my research, who volunteered their precious time and gave good feedback.

  Furthermore, I would like to thank my friends who stood by me, no matter where they are in the world.
  I am particularly indebted to Chihiro for her great assistance and selfless sacrifice.
  In addition, I am grateful to the Makimuras, as well as the other members of KAHF, who became like a second family to me, a home away from home.

  I also thank the Kiyomizudera Temple for awarding me a scholarship, which helped reduce the burden.

  Finally, I would like to express my sincere gratitude to my loved ones, who were with me every step of the way.
  This accomplishment would not have been possible without them.

  There are many more people who deserve thanks. 
  So to everyone who helped support me, knowingly and unknowingly, directly and indirectly, you have my appreciation.

\end{abstract}
\vfill\vfill\null


% Table of contents etc
\clearpage
\tableofcontents
\clearpage
\listoffigures
\listoftables

% Preface
\glsresetall  % Reset acronym usage.
